Метод главных компонент во всех библиотечных реализациях называется РСА, существует для понижения размерности.
Вспомним зачем вообще интересно понижать размерность

\subsection{Связь с SVD}
Причины понижения размерности
\begin{enumerate}
    \item В задачах ML размерность данных может быть очень большой.
    Например, миллионы признаков, особенно когда это разреженные признаки -- становится сложно работать.
    \item Такие данные сложно визуализировать, сложно обучать модели.
    \item Чтобы снизить размерность, нам требуется использовать какой-то механизм, который по-хорошему воспроизводим.
\end{enumerate}
РСА -- линейный способ понижения размерности, который интуитивен и используется широко на практике.
Зачем нам может понадобиться матричное разложение на практике?
Матрица объект-признак $X$ может быть приближена произведением двух матриц меньшего ранга
\[ X_{l,d} = U_{l, k} V_{k, d}^T \]
Здесь размерность $k < d, l$.
В таком случае мы можем сказать, что у нас есть матрица перехода и матрица, которая кодирует всю информацию.
И теперь использовать приближение меньшего ранга для всех наших данных, потеряв часть информации, но используя
меньше признаков.
Такое разложение можно построить с помощью минимизации $\abs{X - U V^T} \rightarrow \min$.
Тут норма Фробениуса или какая-нибудь ещё.
РСА нам говорит, что мы можем матрицу признаков разложить на произведение трёх матриц: $X = U \Sigma V^T$
В этом выражении матрицы $U, V^T$ являются ортогональными матрицами, а $\Sigma $ -- диагональная матрица.
Можем сказать, что $k$ ненулевых признаков остаются ненулевыми и значимыми в ходе такого SVD разложения.
Получаем это с помощью взятия $k$ диагональных элементов $X$ и соответствующих им столбцов в матрицах $U, V$.
Таким образом при перемножении получается итоговая матрица меньшей размерности, итоговый ранг равен $k$

\subsection{Теорема Эккарта–Юнга, доказательство}

\underline{Постановка задачи}

Рассмотрим матрицу $\bm{X} \in \mathbb{R}^{n \times d}$ с центрированными столбцами (нулевым средним).
Мы хотим найти матрицу $\hat{\bm{X}} \in \mathbb{R}^{n \times d}$ ранга $k < \text{rank}(\bm{X})$ такую, что:
\[ \hat{\bm{X}} = \arg\min_{\text{rank}(\bm{B}) \leq k} \|\bm{X} - \bm{B}\|_F^2\]
Здесь $\|\cdot\|_F$ --- норма Фробениуса: $\|\bm{A}\|_F = \sqrt{\sum_{i,j} a_{ij}^2}$.

\underline{Формулировка теоремы}

\begin{theorem}[Эккарта-Юнга, 1936]
    Пусть $\bm{X} \in \mathbb{R}^{n \times d}$ имеет сингулярное разложение:
    \[ \bm{X} = \bm{U}\bm{\Sigma}\bm{V}^\top = \sum_{i=1}^{r} \sigma_i \bm{u}_i \bm{v}_i^\top \]
    где $r = \text{rank}(\bm{X})$, $\sigma_1 \geq \sigma_2 \geq \dots \geq \sigma_r > 0$ --- сингулярные числа,
    $\bm{U} = [\bm{u}_1, \dots, \bm{u}_n]$, $\bm{V} = [\bm{v}_1, \dots, \bm{v}_d]$ --- ортогональные матрицы.
    
    Тогда для любого $k < r$ оптимальная аппроксимация ранга $k$ в смысле нормы Фробениуса задаётся как:
    \[ \hat{\bm{X}}_k = \sum_{i=1}^{k} \sigma_i \bm{u}_i \bm{v}_i^\top = \bm{U}_k \bm{\Sigma}_k \bm{V}_k^\top \]
    Здесь $\bm{U}_k$, $\bm{V}_k$ --- первые $k$ столбцов $\bm{U}$ и $\bm{V}$, $\bm{\Sigma}_k = \text{diag}(\sigma_1, \dots, \sigma_k)$.
    
    При этом минимальная ошибка равна:
    \[ \|\bm{X} - \hat{\bm{X}}_k\|_F^2 = \sum_{i=k+1}^{r} \sigma_i^2 \]
\end{theorem}

\subsubsection{способ через собственное разложение или через подстановку}
\underline{Доказательство способом через собственное разложение}

\begin{proof}
    Рассмотрим два подхода к доказательству.
    
    \underline{Подход 1: Через спектральное разложение $\bm{X}^\top\bm{X}$}
    
    Пусть $\bm{X} = \bm{U}\bm{\Sigma}\bm{V}^\top$ --- SVD\@.
    Тогда:
    \[ \bm{X}^\top\bm{X} = \bm{V}\bm{\Sigma}^2\bm{V}^\top \]
    где $\bm{\Sigma}^2 = \text{diag}(\sigma_1^2, \dots, \sigma_r^2, 0, \dots, 0)$.
    
    Для любой матрицы $\bm{B}$ ранга не более $k$:
    \[ \|\bm{X} - \bm{B}\|_F^2 = \text{tr}((\bm{X} - \bm{B})^\top(\bm{X} - \bm{B})) \]
    
    Зафиксируем $\bm{B}$ и рассмотрим её SVD: $\bm{B} = \bm{P}\bm{\Lambda}\bm{Q}^\top$, где $\bm{P} \in \mathbb{R}^{n \times k}$, $\bm{Q} \in \mathbb{R}^{d \times k}$ имеют ортонормированные столбцы, $\bm{\Lambda} = \text{diag}(\lambda_1, \dots, \lambda_k)$.
    
    Тогда
    \begin{align*}
        \|\bm{X} - \bm{B}\|_F^2 &= \|\bm{X}\|_F^2 + \|\bm{B}\|_F^2 - 2\text{tr}(\bm{X}^\top\bm{B}) \\
        &= \sum_{i=1}^r \sigma_i^2 + \sum_{i=1}^k \lambda_i^2 - 2\text{tr}(\bm{X}^\top\bm{P}\bm{\Lambda}\bm{Q}^\top)
    \end{align*}
    
    По неравенству Коши-Шварца для матриц
    \[ \text{tr}(\bm{X}^\top\bm{P}\bm{\Lambda}\bm{Q}^\top) \leq \sum_{i=1}^k \sigma_i \lambda_i \]
    Здесь равенство достигается при $\bm{P} = \bm{U}_k$, $\bm{Q} = \bm{V}_k$, $\bm{\Lambda} = \bm{\Sigma}_k$.
    
    Таким образом:
    \[ \min_{\text{rank}(\bm{B}) \leq k} \|\bm{X} - \bm{B}\|_F^2 \geq \sum_{i=1}^r \sigma_i^2 - \sum_{i=1}^k
    \sigma_i^2 = \sum_{i=k+1}^r \sigma_i^2 \]
    
    И эта нижняя граница достигается при $\bm{B} = \bm{U}_k\bm{\Sigma}_k\bm{V}_k^\top$.
    
    \underline{Подход 2: Через подпространства.}
    
    Рассмотрим проекцию $\bm{X}$ на подпространство, натянутое на первые $k$ правых сингулярных векторов:
    \[ \hat{\bm{X}}_k = \bm{X}\bm{V}_k\bm{V}_k^\top \]
    
    Покажем, что это минимизирует ошибку.
    Для любой матрицы $\bm{B}$ ранга $k$ пусть $\mathcal{V}$ --- подпространство, натянутое на её правые сингулярные
    векторы.
    Тогда:
    \[ \|\bm{X} - \bm{B}\|_F^2 \geq \|\bm{X} - \bm{X}\bm{P}_{\mathcal{V}}\|_F^2 \]
    Здесь $\bm{P}_{\mathcal{V}}$ --- проектор на $\mathcal{V}$.
    
    Но $\|\bm{X} - \bm{X}\bm{P}_{\mathcal{V}}\|_F^2 = \|\bm{X}\|_F^2 - \|\bm{X}\bm{P}_{\mathcal{V}}\|_F^2$, и:
    \[ \|\bm{X}\bm{P}_{\mathcal{V}}\|_F^2 = \text{tr}(\bm{P}_{\mathcal{V}}\bm{X}^\top\bm{X}\bm{P}_{\mathcal{V}}) =
    \text{tr}(\bm{P}_{\mathcal{V}}\bm{V}\bm{\Sigma}^2\bm{V}^\top\bm{P}_{\mathcal{V}}) \]
    
    Максимум этого выражения (и, соответсвенно, минимум нормы исходной разницы) достигается, когда
    $\mathcal{V}$ --- подпространство, натянутое на первые $k$ правых сингулярных векторов $\bm{v}_1, \dots, \bm{v}_k$.
    
    Таким образом, $\hat{\bm{X}}_k = \bm{X}\bm{V}_k\bm{V}_k^\top$ является оптимальной аппроксимацией.
\end{proof}

\underline{Следствия теоремы}

\begin{enumerate}
    \item \textbf{Оптимальность PCA:} PCA находит оптимальное (в смысле нормы Фробениуса) $k$-мерное подпространство
    для аппроксимации данных.
    \item \textbf{Доля объяснённой дисперсии:} доля дисперсии, сохранённая при использовании $k$ главных компонент:
    \[ R_k^2 = \frac{\sum_{i=1}^k \sigma_i^2}{\sum_{i=1}^r \sigma_i^2}\]
    \item \textbf{Вложенная оптимальность:} для любого $k_1 < k_2$, оптимальная аппроксимация ранга $k_2$ содержит
    в себе оптимальную аппроксимацию ранга $k_1$.
\end{enumerate}

\subsection{Две альтернативные итеративные формулировки}

\underline{Формулировка 1: Последовательная максимизация дисперсии}

Алгоритм может быть сформулирован как последовательное нахождение главных компонент:

\begin{enumerate}
    \item Первая главная компонента:
    \[ \bm{v}_1 = \arg\max_{\|\bm{v}\|=1} \|\bm{X}\bm{v}\|^2 = \arg\max_{\|\bm{v}\|=1} \bm{v}^\top\bm{X}^\top\bm{X}\bm{v} \]
    
    \item Для $k$-й компоненты ($k > 1$):
    \[ \bm{v}_k = \arg\max_{\|\bm{v}\|=1, \bm{v} \perp \{\bm{v}_1, \dots, \bm{v}_{k-1}\}} \|\bm{X}\bm{v}\|^2 \]
\end{enumerate}

\underline{Формулировка 2: Минимизация ошибки реконструкции}

Альтернативно, PCA можно сформулировать как задачу минимизации при условии $\bm{W}^\top\bm{W} = \bm{I}_k$:

\[ \min_{\bm{W} \in \mathbb{R}^{d \times k}, \bm{Z} \in \mathbb{R}^{n \times k}} \|\bm{X} - \bm{Z}\bm{W}^\top\|_F^2 \]

Здесь:
\begin{itemize}
    \item $\bm{W}$ --- матрица нагрузок (loadings)
    \item $\bm{Z}$ --- матрица счётных факторных значений (scores)
\end{itemize}

\underline{Связь между формулировками}

Две формулировки эквивалентны.
Из решения задачи минимизации ошибки реконструкции:
\[ \bm{Z} = \bm{X}\bm{W}, \quad \bm{W} = \bm{V}_k \]

\underline{Итеративный алгоритм NIPALS}

Алгоритм NIPALS (Nonlinear Iterative Partial Least Squares) вычисляет главные компоненты итеративно:

\begin{enumerate}
    \item Инициализировать $\bm{t}$ (например, первый столбец $\bm{X}$)
    \item Повторять до сходимости:
    \begin{align*}
        \bm{p} &\leftarrow \frac{\bm{X}^\top\bm{t}}{\bm{t}^\top\bm{t}} \quad \text{(нагрузки)} \\
        \bm{p} &\leftarrow \frac{\bm{p}}{\|\bm{p}\|} \quad \text{(нормализация)} \\
        \bm{t} &\leftarrow \bm{X}\bm{p} \quad \text{(счётные значения).}
    \end{align*}
    \item Вычесть компоненту: $\bm{X} \leftarrow \bm{X} - \bm{t}\bm{p}^\top$
    \item Перейти к следующей компоненте
\end{enumerate}

\subsection{Предобработка для PCA}

\underline{Важность предобработки}

PCA чувствителен к масштабированию признаков.
Без правильной предобработки компоненты могут быть определены преимущественно признаками с наибольшей дисперсией, что не
всегда соответствует их информативности.

\underline{Основные методы предобработки}

\begin{table}[h]
    \centering
    \begin{tabular}{|p{0.25\textwidth}|p{0.35\textwidth}|p{0.3\textwidth}|}
        \hline
        \textbf{Метод}                         & \textbf{Формула}                                           & \textbf{Когда использовать}                \\ \hline
        Центрирование                          & $\tilde{x}_{ij} = x_{ij} - \bar{x}_j$                      & Всегда для PCA                             \\\hline
        Масштабирование (StandardScaler)       & $\tilde{x}_{ij} = \frac{x_{ij} - \bar{x}_j}{s_j}$ & Когда признаки в разных единицах измерения \\\hline
        Масштабирование к [0,1] (MinMaxScaler) & $\tilde{x}_{ij} = \frac{x_{ij} - \min_j}{\max_j - \min_j}$ & Когда важны относительные различия \\\hline
        Нормализация (L2)                      & $\tilde{\bm{x}}_i = \frac{\bm{x}_i}{\|\bm{x}_i\|}$         & Для данных с разной длиной векторов        \\\hline
    \end{tabular}
    \caption{Методы предобработки для PCA}
    \label{tab:methods}
\end{table}

\subsection*{Центрирование: необходимое условие}

Для PCA \textbf{обязательно} требуется центрирование данных.
Без этого первая главная компонента будет направлена к центру масс данных, а не в направлении максимальной дисперсии.

\[ \bar{\bm{x}} = \frac{1}{n} \sum_{i=1}^n \bm{x}_i, \quad \tilde{\bm{x}}_i = \bm{x}_i - \bar{\bm{x}} \]

\underline{Когда масштабировать}

\begin{enumerate}
    \item Признаки измерены в разных единицах (метры vs килограммы)
    \item Дисперсии признаков сильно различаются
    \item Нет априорных причин считать один признак более важным
\end{enumerate}

\underline{Когда не масштабировать:}
\begin{enumerate}
    \item Все признаки измерены в одних единицах
    \item Различия в дисперсии отражают реальную важность признаков.
    \item Хотим сохранить оригинальную шкалу
\end{enumerate}

\underline{Пример: влияние масштабирования}

Рассмотрим два признака:
\begin{itemize}
    \item $x_1$: рост в сантиметрах (дисперсия $\approx 100$)
    \item $x_2$: вес в килограммах (дисперсия $\approx 100$)
\end{itemize}

Без масштабирования PCA будет определяться обоими признаками.
Если изменить единицы измерения:
\begin{itemize}
    \item $x_1$: рост в метрах (дисперсия $\approx 0.01$)
    \item $x_2$: вес в килограммах (дисперсия $\approx 100$)
\end{itemize}

Без масштабирования первая компонента будет почти полностью определяться весом.

\underline{Обработка пропущенных значений}

Перед применением PCA необходимо обработать пропущенные значения
\begin{enumerate}
    \item Удаление объектов с пропусками (при небольшом их количестве)
    \item Заполнение средним/медианой значением
    \item PCA с пропусками (опционально, сложные методы)
\end{enumerate}

\underline{Обработка выбросов}.
PCA чувствителен к выбросам.
Методы обработки
\begin{enumerate}
    \item Использование RobustScaler
    \item Методы robust PCA
    \item Удаление выбросов перед анализом
\end{enumerate}

Перед применением PCA стоит проверить линейные зависимости

\underline{Практический конвейер предобработки}

\begin{enumerate}
    \item Обработка пропущенных значений
    \item Обнаружение и обработка выбросов
    \item Проверка и преобразование типов данных
    \item Кодирование категориальных переменных (если необходимо)
    \item Центрирование данных
    \item Масштабирование (при необходимости)
    \item Проверка корреляций между признаками
\end{enumerate}

\underline{Визуализация необходимости предобработки}

\begin{enumerate}
    \item \textbf{Boxplot:} Для обнаружения различий в масштабах признаков
    \item \textbf{Scatterplot matrix:} Для визуализации отношений между признаками
    \item \textbf{Гистограммы:} Для проверки распределений
    \item \textbf{Корреляционная матрица:} Для выявления мультиколлинеарности
\end{enumerate}